\documentclass[10pt,a4paper]{amsart}
\usepackage[margin=19mm]{geometry}
\usepackage{amsmath,amssymb,mathtools}
\usepackage[scr=rsfs,cal=euler]{mathalfa}
\usepackage{xfrac}
\usepackage[unicode,hidelinks,bookmarks=false]{hyperref}
\newcommand{\ie}{i.e.\ }
\newcommand{\cf}{cf.\ }
\newcommand{\resp}{respectively\ }
\newcommand{\Subsection}{\subsection}
\providecommand{\nobreakdash}{\nobreak\hbox{-}}
\setlength{\emergencystretch}{2em}
\multlinegap0pt
\usepackage{bm}
\usepackage{bbm}
\usepackage{dsfont}
\usepackage{xspace}
\usepackage{cancel}
\usepackage{mathtools}
\usepackage{amsthm}
\makeatletter
\@ifundefined{theorem}{\newtheorem{theorem}{Theorem}}{}
\makeatother
\newcommand{\Uadh}{{U^{\textup{ad}}_h}}
\newcommand{\Uad}{{U^{\textup{ad}}}}
\newcommand{\bW}{{\mathbf{W}}}
\newcommand{\bn}{{\mathbf{n}}}
\newcommand{\bp}{{\mathbf{p}}}
\newcommand{\bq}{{\mathbf{q}}}
\newcommand{\br}{{\mathbf{r}}}
\newcommand{\cE}{{\mathcal{E}}}
\newcommand{\cT}{{\mathcal{T}}}
\newcommand{\hF}{{\frac{1}{2}}}
\newcommand{\set}[2]{\left\{{#1}\,:~{#2}\right\}}
\newcommand{\x}[2][blue]{#2}
\title[A Local Discontinuous Galerkin Method for Dirichlet Boundary]{A Local Discontinuous Galerkin Method for Dirichlet Boundary Control Problems}
\author{Peter Benner, Michael Hinze, Hamdullah Y\"ucel}
\date{Source: arXiv:2601.18969v2 (CC BY 4.0)}
\begin{document}
\maketitle

\noindent\emph{Source and licence.} A Local Discontinuous Galerkin Method for Dirichlet Boundary Control Problems. Version arXiv:2601.18969v2 (CC BY 4.0), distributed under the Creative Commons Attribution 4.0 International licence, as recorded in the arXiv licence metadata for this exact version and shown on the abstract page https://arxiv.org/abs/2601.18969v2. Full rights notice: licenses/arxiv-2601.18969-CC-BY-4.0.txt.\par\smallskip

\noindent\textit{Editorial scope.} Counted unit: the Theorem whose source label is \texttt{thm:fully-discrete} in the article (source0.tex lines 956--981, source label \texttt{thm:fully-discrete}), delivered together with the article's own complete original proof of it (source0.tex lines 982--1064, one \texttt{proof} environment of 83 lines). No mathematics is written, repaired or re-derived: statement and proof are the article's own text line for line. Required context retained uncounted: p13b (lines 237--246); DG1 (lines 320--326); DG8a (lines 489--505); DG8c (lines 489--505); DG8d (lines 489--505); eqn:reduced\_problem (lines 561--563); def:quasi\_interpolation (lines 645--647); thm:nodisc\_adjoint\_representation (lines 797--809). Every definition, display and label of those lines is retained unchanged. Omitted from this package and not redistributed: 1--236, 247--319, 327--488, 506--560, 564--644, 648--796, 810--955, 1065--2292. Nothing inside a retained excerpt is omitted or altered: every retained body is delivered exactly as the article's lines are written, so no omission receipt of any kind accompanies this package. Citations: the retained excerpts carry no citation command at all, so no bibliography entry of the article is transcribed and no reference is redistributed. Reference adaptation: none was needed and none was made; every reference inside the delivered unit resolves inside it, so no unresolved reference or citation is produced. Printed slips kept literal: no printed word, symbol, hypothesis or number of the counted statement or of its proof was altered. Layout and class: the article's own class is \texttt{article} with packages including todonotes, a4wide, hyperref, amsmath, amsthm, amsfonts, amssymb, mathtools, algorithm, algorithmic, epsfig, inputenc, graphics, caption; the wrapper is the lane's pinned \texttt{amsart} editorial wrapper at 10pt on A4, which restates the theorem-like environments the retained text needs and adds the article's own macro definitions that the retained text needs (\texttt{\textbackslash Uad}, \texttt{\textbackslash Uadh}, \texttt{\textbackslash bW}, \texttt{\textbackslash bn}, \texttt{\textbackslash bp}, \texttt{\textbackslash bq}, \texttt{\textbackslash br}, \texttt{\textbackslash cE}, \texttt{\textbackslash cT}, \texttt{\textbackslash hF}, \texttt{\textbackslash set}, \texttt{\textbackslash x}), with extra wrapper packages (bm, bbm, dsfont, xspace, cancel, mathtools). The delivered numbering is the unit's own. Assets: no retained span contains a figure environment, a graphics-inclusion command, a PDF-page inclusion, a table or a TikZ picture; no image file is copied into this package, the package declares no asset dependency and no image file is redistributed with it. Licence: version arXiv:2601.18969v2 of this article is distributed under the Creative Commons Attribution 4.0 International licence, as recorded in the arXiv licence metadata for this exact version and shown on the abstract page https://arxiv.org/abs/2601.18969v2. A full rights notice is delivered at licenses/arxiv-2601.18969-CC-BY-4.0.txt. Characters: every retained excerpt is pure ASCII, so the delivered engine, XeLaTeX with its default Latin Modern text font, reports no missing character.
\begin{align}\label{p13b}
\begin{aligned}
\nabla \cdot (-\beta\, z - \epsilon^{\hF} \mathbf{p}) + \alpha\, z &= y - y^{d}
    &&\text{in }\Omega,\\[2mm]
\mathbf{p} &= \epsilon^{\hF} \nabla z
    &&\text{in }\Omega,\\[2mm]
z &= 0
    &&\text{on }\Gamma .
\end{aligned}
\end{align}

\begin{subequations}\label{DG1}
\begin{align}
    \mathbf{W}_h &= \set{\mathbf w \in \big(L^2(\Omega)\big)^2}{ \mathbf w\mid_{K}\in \big( \mathbb{S}^1(K) \big)^2, \quad \forall K \in \cT_h}, \\[4pt]
    V_h &= \set{v \in L^2(\Omega)}{ v \mid_{K}\in \mathbb{S}^1(K), \quad \forall K \in \cT_h}, \\[4pt]
    U_h &= \set{u \in L^2(\Gamma)}{ u \mid_{E}\in \mathbb{S}^1(E), \quad \forall E \in \cE_h^{\partial}},
\end{align}
\end{subequations}

\begin{align}
a_h(\bq_h, \br) + b_h(y_h, \br) 
&= m_{h,1}(u_h, \br)
&& \forall \br \in \bW_h, \label{DG8a} \\[0.3em] 
-b_h(v, \bq_h) + c_h(y_h, v) 
&= m_{h,2}(u_h, v) + F(v)
&& \forall v \in V_h, \label{DG8b} \\[0.3em] 
a_h(\bp_h, \psi) - b_h(z_h, \psi) 
&= 0 
&& \forall \psi \in \bW_h, \label{DG8c} \\[0.3em] 
b_h(\phi, \bp_h) + c_h(\phi, z_h) 
&= (y_h - y^d, \phi)_\Omega
&& \forall \phi \in V_h, \label{DG8d} \\[0.3em] 
\langle \omega u_h + m'_{1,h}(u_h, \bp_h) + m'_{2,h}(u_h, z_h), \, w - u_h \rangle_\Gamma 
&\ge 0
&& \forall w \in \Uadh. \label{DG8e}
\end{align}

\begin{equation}\label{eqn:reduced_problem}
    \min_{u \in \Uad} \widehat J(u).
\end{equation}

\begin{equation}\label{def:quasi_interpolation}
    \pi_h u = \sum_{n \in \mathcal{N}_h^\partial} \frac{(u, \xi_n)}{(1, \xi_n)} \, \xi_n,
\end{equation}

\begin{theorem}\label{thm:nodisc_adjoint_representation}
Let $u$ be the solution of~\eqref{eqn:reduced_problem}, and let 
$u_h$ be the corresponding discrete control obtained with the choice $U_h = U = L^2(\Gamma)$. 
Moreover, let $y$ and $y_h$ denote the corresponding continuous and discrete state solutions, respectively. 
Then, 
\begin{equation}\label{vge}
    \omega \| u - u_h \|_{0,\Gamma}^2  + \|y-y_h\|_{0,\Omega}^2 \le \|y-y_h(u)\|_{0,\Omega}^2 + C(\omega) 
     \big( \epsilon \|(\bp_h(u) - \bp)\cdot\bn\|_{0,\Gamma}^2 
         + \kappa_z^2\|z_h(u) - z\|_{0,\Gamma}^2 \big),
\end{equation}
where the positive constant $C$ is independent of the mesh size~$h$, and where $\bp_h(u):=\bp_h(y(u),\bq(u))$ and $z_h(u):=z_h(y(u),\bq(u))$ are local discontinuous Galerkin approximations of $\bp$ and $z$ associated with the continuous data $u, y(u)$, and $\bq(u)$, respectively, i.e., $\big(z_h(u),\bp_h(u)\big)$ solves \eqref{DG8c}-\eqref{DG8d} with right-hand side $y-y^d$.

\end{theorem}

\begin{theorem}\label{thm:fully-discrete}
Let $u$ be the solution of~\eqref{eqn:reduced_problem}, and let
$u_h$ be the corresponding discrete control obtained with the choice of $U_h$ defined in~\eqref{DG1}.
Further, let $y$ and $y_h$ denote the corresponding \x{continuous and discrete states}, respectively. 
Then, 
\begin{multline}\label{gge}
    \omega \| u - u_h \|_{0,\Gamma}^2  + \|y-y_h\|_{0,\Omega}^2 \le \|y-y_h(u)\|_{0,\Omega}^2 + C(\omega) 
     \big( \epsilon\|(\bp_h(u) - \bp)\cdot \bn\|_{0,\Gamma}^2 
         + \kappa_z^2\|z_h(u) - z\|_{0,\Gamma}^2 \big)  \\ + C(\omega) \|\pi_hu-u\|_{0,\Gamma}^2 + \big(\epsilon^{1/2}\|(\bp_h - \bp(y_h))\cdot \bn\|_{0,\Gamma}+ \kappa_z\|z_h - z(y_h)\|_{0,\Gamma}\big)\|\pi_hu-u\|_{0,\Gamma}  \\ + C \big(\|u\|_{s,\Gamma}, \epsilon^{1/2}\|\bp(u)\cdot\bn\|_{s,\Gamma}\big)\,
       \|\pi_h u - u\|_{-s,\Gamma},
\end{multline}
% the following estimate holds:
% \begin{align}
% \|u - u_h\|_{0, \Gamma} + 
% \|y - y_h\|_{0,\Omega}
%   &\leq 
%      C\Big(
%         \|y_h(u) - y\|_{0,\Omega}
%         + \|\bp_h(u) - \bp\|_{0,\Gamma}
%         + \|z_h(u) - z\|_{0, \Gamma}
%       \Big)  \nonumber\\
%   &\quad
%      + C\, h^{s}\,\|u\|_{s,\Gamma},
% \end{align}
where the triplet $\big(y_h(u), z_h(u),\bp_h(u)\big)$ denotes the local discontinuous Galerkin approximations of $(y,z,\bp)$ for a given $u\in H^s(\Gamma)$  with $0\le s\le 1$, i.e., $\big(y_h(u), z_h(u),\bp_h(u)\big)$ solves \eqref{DG8a}-\eqref{DG8d} with the data $u$ and $y-y^d$, and $(\bp(y_h),z(y_h))$ denotes the continuous adjoint flux and potential associated with the optimal discrete state $y_h$, i.e., the solution to \eqref{p13b} with right-hand side $y_h-y^d$.
\end{theorem}

\begin{proof}
From the definitions of the variational inequalities and the quasi-interpolation property in \eqref{def:quasi_interpolation}, we have 
\begin{align*}
 \widehat J'(u)(u_h - u) &\ge 0 \qquad \forall\, u_h \in \Uadh = U_h \cap \Uad \subset \Uad, \\
 \widehat J_h'(u_h)(\pi_h u - u_h) &\ge 0 \qquad \forall\, \pi_h u \in \Uadh.
\end{align*}
Adding these inequalities gives
\begin{equation}\label{eqn:f1}
   \underbrace{\big(\widehat J_h'(u_h) - \widehat J'(u)\big)(u - u_h)}_{(1)} + \underbrace{\widehat J_h'(u_h)(\pi_h u - u)}_{(2)} \ge 0.
\end{equation}
We now can estimate (1) in \eqref{eqn:f1} as in the proof of Theorem \ref{thm:nodisc_adjoint_representation}, which gives contributions to the upper part of \eqref{gge}. %Using the identity
% \[
%  \big(\widehat J_h'(u_h) - J'(u)\big)(u - u_h) 
%  = \omega \langle u_h - u,\, u - u_h \rangle_{\Gamma} 
%    + m_{h,1}(u - u_h,\, \bp_h - \bp) 
%    + m_{h,2}(u - u_h,\, z_h - z),
% \]
% and substituting into \eqref{eqn:f1}, we obtain
% \begin{align}\label{eqn:f2}
% \omega \|u - u_h\|_{\Gamma}^2 
% &\le m_{h,1}(u - u_h,\, \bp_h - \bp) 
%    + m_{h,2}(u - u_h,\, z_h - z)
%    + \widehat J_h'(u_h)(\pi_h u - u) \nonumber \\
% &= m_{h,1}(u - u_h,\, \bp_h - \bp_h(u)) 
%    + m_{h,2}(u - u_h,\, z_h - z_h(u)) \nonumber\\
% &\quad + m_{h,1}(u - u_h,\, \bp_h(u) - \bp) 
%    + m_{h,2}(u - u_h,\, z_h(u) - z)
%    + \widehat J_h'(u_h)(\pi_h u - u) \nonumber \\
% &= I_1 + I_2 + I_3,
% \end{align}
% where $\bp_h(u):=\bp_h(y(u),\bq(u))$ and $z_h(u):=z_h(y(u),\bq(u))$ are local discontinuous Galerkin approximations of $\bp$ and $z$, respectively.\\
% \noindent To derive a bound for $I_1$ in \eqref{eqn:f2}, we make use of the optimality conditions  together with Young's inequality
% \begin{align}\label{eqn:f3}
% I_1 
% &= m_{h,1}(u - u_h,\, \bp_h - \bp_h(u)) 
%    + m_{h,2}(u - u_h,\, z_h - z_h(u)) \nonumber \\[4pt]
% &= a_h(\bq_h(u) - \bq_h,\, \bp_h + \bp_h(u)) 
%    + b_h(y_h(u) - y_h,\, \bp_h - \bp_h(u)) \nonumber \\
% &\quad - b_h(z_h - z_h(u),\, \bq_h(u) - \bq_h) 
%    + c_h(y_h(u) - y_h,\, z_h - z_h(u)) \nonumber \\[4pt]
% &= (y_h - y,\, y_h(u) - y + y - y_h)_{\Omega} \nonumber \\[4pt]
% &\le -\|y - y_h\|^2_{0,\Omega} 
%    + \frac{1}{2}\|y - y_h\|^2_{0,\Omega} 
%    + \frac{1}{2}\|y_h(u) - y\|^2_{0,\Omega},
% \end{align}
% where $y_h(u)$ denotes the local discontinuous Galerkin approximation of the state variable corresponding to the Dirichlet boundary condition $u \in H^{s}(\Gamma)$. \\
% \noindent From the definitions of the (bi)linear forms and Young's inequality, the term $I_2$ in \eqref{eqn:f2} can be estimated as
% \begin{align}\label{eqn:f4}
%  I_2 
%  &= m_{h,1}(u - u_h,\, \bp_h(u) - \bp) 
%    + m_{h,2}(u - u_h,\, z_h(u) - z) \nonumber \\[0.3em]
%  &\le \frac{\omega}{4}\,\|u - u_h\|_{\Gamma}^2 
%    + C(\omega, \epsilon, \kappa_z)\,
%      \big( \|\bp_h(u) - \bp\|_{\Gamma}^2 
%           + \|z_h(u) - z\|_{\Gamma}^2 \big).
% \end{align}
To treat (2) in \eqref{eqn:f1} we write
\begin{align}\label{eqn:f5}
J_h'(u_h)(\pi_h u - u)
  &= \langle\, \omega u_h - \epsilon^{\hF} \, \bp_h \cdot \bn + \kappa_z z_h ,\, \pi_h u - u \,\rangle \nonumber \\[0.3em]
  &= \langle\, \omega (u_h-u),\, \pi_h u - u \,\rangle \nonumber\\
  &\quad
     + \langle\, \epsilon^{\hF} \, (\bp(y_h) -\bp_h)\cdot \bn 
       + \kappa_z (z_h - z(y_h)),\, \pi_h u - u \,\rangle  \nonumber\\
  &\quad
     + \langle\, \epsilon^{\hF} \, (\bp(u)-\bp(y_h))\cdot \bn
       + \kappa_z (z(y_h) - z(u)),\, \pi_h u - u \,\rangle  \nonumber\\
  &\quad
     + \langle\, \omega u -  \epsilon^{\hF} \,\bp(u)\cdot \bn + \kappa_z z(u),\, \pi_h u - u \,\rangle.
\end{align}
Here, $\bp(u),z(u)$ denote the optimal adjoint flux and adjoint potential, respectively, while $\bp(y_h),z(y_h)$ denote auxiliary continuous adjoint flux and potential associated with the optimal discrete adjoint state $y_h$, i.e., with right-hand side $y_h - y^d \in L^2(\Omega)$ in \eqref{p13b}. We already  note at this stage of the proof that the continuous adjoint variable $z$ is zero on the boundary. Straightforward estimation gives
\begin{multline}\label{extra-term}
J_h'(u_h)(\pi_h u - u) \le \frac{\omega}{4}\|u-u_h\|_{0,\Gamma}^2+C(\omega) \|\pi_hu-u\|_{0,\Gamma}^2  \\ +C \big(\epsilon^{1/2}\|(\bp_h - \bp(y_h))\cdot\bn\|_{0,\Gamma}+\kappa_z\|z_h - z(y_h)\|_{0,\Gamma}\big)\|\pi_hu-u\|_{0,\Gamma}  \\ + C\epsilon^{1/2}\|(\bp(y_h)-\bp(u))\cdot\bn\|_{0,\Gamma}\|\pi_hu-u\|_{0,\Gamma}  \\ + C\|(u,\epsilon^{1/2}\bp(u)\cdot\bn)\|_{s,\Gamma}\|\pi_hu-u\|_{-s,\Gamma}.
\end{multline}
Since $y_h, y^d\in L^2(\Omega)$, we have 
$z(y_h) \in H^2(\Omega) \cap H_0^1(\Omega)$. The trace theorem together with the continuity of the continuous adjoint solution operator gives 
\begin{equation}\label{eqn:f7}
\|(\bp(y_h) - \bp(u))\cdot\bn\|_{0,\Gamma}\leq  C \|z(y_h) - z(u)\|_{2,\Omega}  \leq C \|y_h - y \|_{0,\Omega},
\end{equation}
so that
\[C \epsilon^{1/2}\|(\bp(y_h) - \bp(u))\cdot\bn\|_{0,\Gamma}\|\pi_hu-u\|_{0,\Gamma} \le \frac{1}{4}\|y-y_h\|_{0,\Omega}^2 + C\|\pi_hu-u\|_{0,\Gamma}^2.\]
Sorting the addends of \eqref{extra-term} in accordingly gives the estimate \eqref{gge}.
\end{proof}

\end{document}
